\documentclass[12pt]{article}
\usepackage[T1]{fontenc}
\usepackage[utf8]{inputenc}
\usepackage{lmodern}
\usepackage{textcomp}
\usepackage{enumitem}
\usepackage{geometry}
\usepackage{graphicx}
\usepackage{amsmath, amssymb}
\geometry{margin=1in}
\begin{document}
\begin{center}
Engineering - Undergraduate Level Assessment 6 \
Total Marks: 40 \
Answer all questions.
\end{center}

\section*{Multiple Choice (10 marks)}
\begin{enumerate}[label=\arabic*.]
\item If a body is cooled below ambient temperature by a system, it can then be used to freeze water. Calculate the maximum air temperature for which freezing takes place.
\begin{enumerate}[label=(\alph*)]
    \item 190^\ensuremath{\circF}
    \item 200^\ensuremath{\circF}
    \item 170^\ensuremath{\circF}
    \item 180^\ensuremath{\circF}
    \item 120^\ensuremath{\circF}
\end{enumerate}
\item A Binary number system has how many digits.
\begin{enumerate}[label=(\alph*)]
    \item 0
    \item 1
    \item 2
    \item 10
\end{enumerate}
\item The energy stored in the magnetic field in a solenoid 30 cm long and 3 cm diameter wound with 1000 turns of wire carrying a current at 10 amp, is
\begin{enumerate}[label=(\alph*)]
    \item 0.015 joule.
    \item 15 joule.
    \item 0.0015 joule.
    \item 1.5 joule.
    \item 0.15 millijoule.
\end{enumerate}
\item If six messages having probabilities (1 / 4), (1 / 4), (1 / 8), (1 / 8), (1 / 8) and (1 / 8) can be accommodated in a single communication system, determine the entropy.
\begin{enumerate}[label=(\alph*)]
    \item 2.00 bits/message
    \item 2.25 bits/message
    \item 1.75 bits/message
    \item 3.00 bits/message
    \item 2.75 bits/message
\end{enumerate}
\item Find f\_1(t) \{\_\ensuremath{\ast}\} f\_2(t) if f\_1(t) = 2 e^-4 tu(t) and f\_2(t) = 5 cos 3 t u(t).
\begin{enumerate}[label=(\alph*)]
    \item f\_1(t) \{\_st\} f\_2(t) = (1.6 cos 3 t + 1.2 sin 3 t - 1.6 e^-4 t)u(t)
    \item f\_1(t) \{\_st\} f\_2(t) = (1.2 cos 3 t + 1.6 sin 3 t - 1.3 e^-4 t)u(t)
    \item f\_1(t) \{\_st\} f\_2(t) = (1.4 cos 3 t + 1.2 sin 3 t - 1.5 e^-4 t)u(t)
    \item f\_1(t) \{\_st\} f\_2(t) = (1.6 cos 3 t + 1.0 sin 3 t - 1.6 e^-4 t)u(t)
\end{enumerate}
\end{enumerate}

\section*{True / False (10 marks)}
\textit{Answer True or False}
\begin{enumerate}[label=\arabic*.]
\setcounter{enumi}{5}
\item True or False: The attenuation per meter along the waveguide for a wavelength of 2 meters is 545 db/meter.
\item True or False: The z-transform of f[n] = 3δ[n+2] + 4δ[n-1] + 2δ[n-3] is F(z) = 4 $z^2$ + 3 z^-1 + 2 z^-3.
\item True or False: The partial-fraction expansion of F(s) is F(s) = [(1 / 4) / (s + 1)] + [(1 / 2) / (s + $1)^2$] - [(1 / 4) / (s + 3)].
\item True or False: A box that shows the output results of a control subsystem is referred to as a process box.
\item True or False: The average energy density of a plane wave with a Poynting vector of 5 watts/$meter^2$ is (1/4) erg/$meter^3$.
\end{enumerate}

\section*{Long Form Response (20 marks)}
\textit{Answer each question with a clear and complete explanation.}
\begin{enumerate}[label=\arabic*.]
\setcounter{enumi}{10}
\item Discuss how to calculate the optimal length of a half-wave dipole antenna for a frequency of 300 MHz, considering different velocity factors. Include the implications of these factors on antenna performance.
\item Calculate the number of 300-lumen lamps needed to achieve an average illuminance of 50 lux in a 4 m x 3 m room. Discuss your calculations and the underlying principles of illuminance.
\end{enumerate}

\vfill
\noindent\textit{End of Paper}
\end{document}